\chapter{The incident detail page and its timeline}
\label{ch:detail}

\shot{06-incident-detail.png}{One incident, open. The timeline runs down the
page; the controls for changing it sit alongside.}

This is the most important screen in the application. Everything else exists to
get you here or to summarise what happened here.

\section{The timeline is the product}

Everything that has ever happened to this incident appears as an entry: it was
opened, someone acknowledged it, severity was raised, a note was posted, it was
mitigated, it was resolved.

\where{web/src/components/Timeline.tsx}
\begin{lstlisting}[language=tsx]
{events.map((event) => {
  const { relative, absolute } = formatTimestamp(event.occurred_at);
  return (
    <li key={event.id} className="relative">
      <span className={dotColour(event.type)} />
      <p className="text-sm text-slate-800 dark:text-slate-200">{event.summary}</p>
      <time dateTime={event.occurred_at ?? undefined}>{relative}</time>
      {typeof event.payload.note === 'string' && event.payload.note.length > 0 && (
        <p>{event.payload.note}</p>
      )}
      {event.request_id && (
        <span>{event.request_id}</span>
      )}
    </li>
  );
})}
\end{lstlisting}

Two details repay attention.

\begin{description}[leftmargin=1.2cm,style=nextline]
\item[Both time formats] \fpath{relative} shows ``12 minutes ago'', which is what
  you want during an incident. \fpath{dateTime} carries the exact machine
  timestamp, which is what you want when writing the postmortem. The
  \fpath{<time>} element is the HTML built for exactly this.
\item[\fpath{request\_id} on the entry] Every timeline event records the ID of
  the HTTP request that created it. When something looks wrong six months later,
  that string leads to the precise server log line. Very few systems do this,
  and it costs one column.
\end{description}

\section{Append-only, enforced three times}

Chapter~\ref{ch:arch} claimed the timeline cannot be edited. Here is the
enforcement.

\where{api/app/Models/IncidentEvent.php}
\begin{lstlisting}[language=PHP]
/**
 * Corrections are appended as new events, never applied in place.
 */
protected static function booted(): void
{
    static::updating(static function (): never {
        throw new LogicException('Incident timeline events are append-only and cannot be modified.');
    });

    static::deleting(static function (): never {
        throw new LogicException('Incident timeline events are append-only and cannot be deleted.');
    });
}
\end{lstlisting}

\fpath{booted()} runs once when the class is first used, and registers two
listeners. Any attempt to save a change to an existing event, or to delete one,
throws --- from anywhere in the codebase, including a future controller written
by someone who never read this book.

\note{The return type \fpath{: never} is doing real work. It tells PHP --- and
any static analyser --- that this function cannot return normally; it always
throws. A reader skimming the file learns the guarantee from the signature
without reading the body.}

\subsection{Why three layers and not one}

\begin{center}
\begin{tabular}{@{}>{\raggedright\arraybackslash}p{3.0cm}>{\raggedright\arraybackslash}p{4.4cm}p{6.4cm}@{}}
\toprule
\textbf{Layer} & \textbf{What it stops} & \textbf{What it cannot stop} \\
\midrule
API & Requests that ask to edit history --- there is no such endpoint. & Code inside the application calling the model directly. \\
Model & Any PHP code in this project, deliberate or accidental. & Someone connecting to PostgreSQL and running \fpath{UPDATE}. \\
Schema & Structure: there is no ``edited\_at'', no soft-delete column, nowhere for a revision to live. & A determined administrator with database credentials. \\
\bottomrule
\end{tabular}
\end{center}

\warnbox{No layer is complete on its own, and the last row is honest: a person
with direct database access can change anything. The guarantee is not
mathematical --- it is that no \emph{ordinary} path through the system can
rewrite history, and that doing so requires deliberate, logged, out-of-band
action rather than a stray line of application code.}

\section{Changing an incident}

The controls beside the timeline are gated by permission, the same way the
``Report incident'' button was in Chapter~\ref{ch:incidents}. But status changes
carry a second gate the list page did not have: even a user \emph{with}
permission may only make a \emph{legal} move.

Recall the transition table from Chapter~\ref{ch:arch}: an incident that is
\fpath{open} may become \fpath{acknowledged} or \fpath{resolved} --- and nothing
else. The frontend shows only those two options. The API independently checks
the same rule and returns \fpath{409 Conflict} for anything else.

\note{Two independent checks of one rule sounds redundant, and it is --- on
purpose. They fail differently. The frontend check fails \emph{invisibly}: the
option is not offered, so the mistake is never made. The API check fails
\emph{loudly}: a 409 with a message. One is for people, the other is for
programs, and only the second is a security control.}

\section{The page updates while you watch it}

Open this incident in two browsers, change the status in one, and the other
updates within a second --- no refresh.

That is the \fpath{realtime} container. The chain, end to end:

\begin{enumerate}[leftmargin=1.5cm]
\item The API commits the change to PostgreSQL.
\item After the commit, it publishes a small event to Redis.
\item \fpath{realtime} is subscribed to Redis and receives it.
\item It writes the event down every open connection belonging to that
  organisation.
\item Each browser receives it and updates the page in place.
\end{enumerate}

\warnbox{Step 4 says ``belonging to that organisation'' and that phrase is
load-bearing. Every connection is authenticated, and events are filtered by
organisation before being written. Without that filter, a live stream is a
spectacular data leak: every customer would receive every other customer's
incidents in real time, and no page in the UI would ever reveal it.}

\subsection{Why the browser cannot just use its normal token}

The realtime connection is opened by the browser's \fpath{EventSource}, which
cannot set an \fpath{Authorization} header. So the frontend first asks the API
for a short-lived \emph{ticket}, then opens the stream with that ticket in the
query string.

\note{A ticket is used rather than the access token because query strings are
not private: they appear in server logs, proxy logs and browser history. A
ticket is valid for seconds, for one connection, and grants nothing but the
right to listen. Leaking it in a log is uninteresting; leaking your access token
there is not.}

\section{Check your understanding}

\begin{itemize}[leftmargin=1.4cm]
\cbox Name the three layers that stop the timeline being edited, and what each
  one misses.
\cbox Why is the same transition rule checked in both the browser and the API?
\cbox Why does a timeline entry store a \fpath{request\_id}?
\cbox Why does the live stream use a ticket instead of the access token?
\end{itemize}
